\section{დასკვნა}

Peared-ის შექმნამ ერთ პროექტში მოაქცია ის, რასაც ჩვეულებრივ რამდენიმე
სხვადასხვა კურსი ფარავს: დისტრიბუციული სისტემები (CRDT,
State Vector Exchange სისტემა, Garbage collection სისტემა), მონაცემთა სტრუქტურები
(აუგმენტირებული B+ ხე), სისტემური პროგრამირება Rust-ზე, ქსელური
პროგრამირება Go-ზე და თანამედროვე დესკტოპ/ვებ ინტერფეისები. საბოლოო
პროდუქტი, ერთობლივი რედაქტორი, რომელიც ცენტრალური სერვისის გარეშე
მუშაობს როგორც ლოკალურ, ისე გლობალურ ქსელში, თავიდანვე დასახულ მიზანს
პასუხობს და რეალურ სცენარში (აუდიტორია არასტაბილური ინტერნეტით) ნამდვილად
გამოსადეგია.

ტექნიკური ცოდნის მიღმა პროექტმა რამდენიმე უფრო ზოგადი გაკვეთილი
გვასწავლა. პირველი, კორექტულობის ინფრასტრუქტურის ღირებულება: debug-
ორაკულმა, protocol drift ტესტებმა და მკაცრმა CI-მ განაწილებული სისტემის
ყველაზე საშიში კლასის შეცდომები (რეპლიკების უხმო დაშორება) განვითარების
ეტაპზევე დაიჭირა, სანამ ისინი „იშვიათ, გაუმეორებელ ბაგებად“ იქცეოდნენ.
მეორე, აბსტრაქციის საზღვრების სწორად გავლება: ის, რომ \texttt{crdt-core}
არაფერს იცის ქსელზე, გეითვეი არაფერს იცის დოკუმენტზე, ხოლო ფრონტენდი
არაფერს იცის CRDT-ზე, სამივე კომპონენტის დამოუკიდებლად ტესტირებას და
პარალელურად განვითარებას აძლევდა საშუალებას. მესამე, გუნდური მუშაობის დისციპლინა: სამი ადამიანის
ნამუშევრის ყოველდღიური შერწყმა pull request-ებით, კოდის განხილვით და
ავტომატური შემოწმებებით თავად პროექტის ისეთივე ნაწილი აღმოჩნდა, როგორც
ალგორითმები.
